\section{Lineage and Context}
\label{sec:receipts-lineage}

\subsection{2016 Language-Model Lesson}
\label{sec:receipts-lineage-lm}

The doctoral lineage of the receipts corpus begins with a 2016 experiment in local language
modelling. The finding of that experiment was precise and consequential: \emph{local text
imitation does not conserve consequence}. A language model trained to predict the next token
in a corpus of enterprise communications can reproduce the surface form of decision language
without preserving the causal chain that makes a decision binding. The model learns to say
``approved'' without knowing what approval requires, what event log would prove approval
occurred, or what downstream process state would change as a result.

This lesson is not merely about language models. It is a general finding about any system
that operates on surface form without consequence binding. The 2016 experiment established
that fidelity of representation is not a substitute for fidelity of consequence. A claim
that looks like an M\&A finding is not an M\&A finding unless it is bound to the event
log that produced it, the process model that defines conformance, and the conformance
algorithm that computed the fitness score. The receipts corpus is the direct institutional
response to this finding. It refuses to allow any claim to enter the board record unless
the consequence chain---from raw event log to signed receipt---is intact and inspectable.

\subsection{Enterprise Process Gap}
\label{sec:receipts-lineage-gap}

The 2016 finding exposed a structural gap between two enterprise realities. The first
reality is that enterprise processes leave event traces: transaction logs, ERP audit
trails, supply-chain timestamps, AR aging records, SLA breach notifications. These traces
contain the ground truth of process behavior. The second reality is that the claims made
in M\&A transactions are typically not derived from those traces. They are derived from
management presentations, normalized EBITDA schedules, and consultant assessments that
may or may not be grounded in the underlying process evidence.

The gap between process evidence and M\&A claims is the target of this research program.
The receipts corpus closes this gap for five specific M\&A claim domains: EBITDA rework
reduction (\$1.25M, fitness $0.982$, precision $0.945$), working capital and AR velocity
(\$1.37M, fitness $0.978$, replication \texttt{PASSED}), SLA penalty exposure (\$450K
cap, late delivery rate $2.14\%$, log coverage $98.4\%$), GRC compliance leakage
(\$0.00, fitness $1.0$, zero violations), and residual process standardization
($97.5\%$ conformance, drift \texttt{STABLE}). Each of these claims was derived from
a process-mining computation over an event log, and each is bound to its evidence by
a signed JSON receipt.

\subsection{Relationship to the Chatman Equation}
\label{sec:receipts-lineage-chatman}

The Chatman Equation $A = \mu(O^*)$ defines the core architectural claim of the doctoral
research: that an autonomous process agent $A$ is properly understood as the expectation
operator $\mu$ applied to the ensemble of all observable process states $O^*$. This equation
places the burden of proof on observability. An agent that claims to produce a result cannot
be treated as having produced that result unless the result is observable---that is, unless
it appears in the event log of the process the agent is supposed to be governing.

The receipts corpus operationalizes the observability requirement of the Chatman Equation.
The BLAKE3 \texttt{target\_log\_hash} field in each JSON receipt binds the claim to a
specific observed event log. The \texttt{process\_model\_hash} binds the conformance
measurement to a specific process model. Together these hashes make the $O^*$ of the
Chatman Equation concrete: it is not an abstract ensemble of process states, but a
specific, content-addressed, tamper-evident record. The fitness score in each receipt
is the empirical realization of $\mu$: the expectation of conformance computed over
the observed log. The Ed25519 signature seals the equation---it asserts that a
specific authority has computed $\mu(O^*)$ and bound the result to an immutable record.

\subsection{Relationship to Receipts and Replay}
\label{sec:receipts-lineage-replay}

Within the full doctoral lineage, ``receipts and replay'' names the second major
contribution after the 2016 language-model lesson. Replay, in the process-mining sense,
refers to token-based replay fitness: the van der Aalst formula
$f = 1 - (\sum m / \sum c) - (\sum r / \sum p)$, where missing tokens and remaining
tokens quantify the gap between an observed trace and a reference process model. A
replay fitness of $1.0$ means every observed trace conforms perfectly to the model;
values below $0.95$ fail the \texttt{ConformanceAdmissionGate}.

The receipts in this corpus are the durable output of replay. Each M\&A claim receipt
carries a \texttt{wasm4pm\_query\_result} that includes the replay fitness value computed
by the conformance authority module. The receipt is the proof that replay occurred, that
it produced a specific numeric result, and that the result was computed over a specific
event log at a specific point in time. Without the receipt, replay is an ephemeral
computation. With the receipt, replay becomes a permanent, auditable fact in the
research program's evidence chain. This relationship---replay produces the number,
the receipt preserves the number with its derivation chain---is the architectural
core of the cryptographic receipt infrastructure documented in this thesis.
